\makeabstract
{
Effects of Reassessment Policy on Emergency Department Performance
}
{
Mayisa Tasnim (1), Emile Ndayishimiye (1), Kristy Gardner (1)
}
{
\textsuperscript{1}Department of Computer Science, Amherst College, Amherst, MA, USA
}
{
Prolonged length of stays (LOS) in emergency departments results in adverse long and short-term effects for high acuity patients, increasing death and "left without being seen" (LWBS) rates. Misdiagnosis, intensified by staff fatigue and patient congestion, is a key factor in high LOS, especially when uncorrected. A proposed solution to reduce LOS and correct misdiagnosis is to reassess patients after post-triage delay. 
}
{
We developed a discrete-event driven simulation of the emergency department (Sort → Registration → Triage → Treatment Zones) with dynamic patient arrivals, variable staffing, and probabilistic misdiagnosis to evaluate reassessment strategies. Three models were tested: (1) baseline with no misdiagnosis or reassessment, (2) misdiagnosis only with increasing error variability and rate, and (3) misdiagnosis with post-triage reassessment at delays from 0 to 240 minutes. Each model was simulated for one year and repeated ten times to measure LOS, LWBS, and deaths.
}
{
Reassessment after 120 minutes optimizes outcomes. When misdiagnosis variability and rate are low to moderate during this delay, deaths and LWBS rates decrease significantly without extending LOS because more high acuity patients in congested zones are corrected. Any shorter delay reduces reassessment efficiency, while longer delays allow more untreated high-acuity patients to deteriorate, increasing LOS due to treatment complications. LWBS rates were less sensitive to reassessment timing.
}
{
Reassessment after 120 minutes in post-triage delay is effective. Future work will compare alternative staffing models for reassessment, such as dedicating staff for reassessment versus rerouting patients back to triage.
}

\newpage
